% \documentclass[journal]{vgtc}                % final (journal style)
\documentclass[review,journal]{vgtc}         % review (journal style)
% \documentclass[review,journal]{vgtc}         % review (journal style)
%\documentclass[widereview]{vgtc}             % wide-spaced review
%\documentclass[preprint,journal]{vgtc}       % preprint (journal style)

%% Uncomment one of the lines above depending on where your paper is
%% in the conference process. ``review'' and ``widereview'' are for review
%% submission, ``preprint'' is for pre-publication, and the final version
%% doesn't use a specific qualifier.

%% Please use one of the ``review'' options in combination with the
%% assigned online id (see below) ONLY if your paper uses a double blind
%% review process. Some conferences, like IEEE Vis and InfoVis, have NOT
%% in the past.

%% Please note that the use of figures other than the optional teaser is not permitted on the first page
%% of the journal version.  Figures should begin on the second page and be
%% in CMYK or Grey scale format, otherwise, colour shifting may occur
%% during the printing process.  Papers submitted with figures other than the optional teaser on the
%% first page will be refused. Also, the teaser figure should only have the
%% width of the abstract as the template enforces it.

%% These few lines make a distinction between latex and pdflatex calls and they
%% bring in essential packages for graphics and font handling.
%% Note that due to the \DeclareGraphicsExtensions{} call it is no longer necessary
%% to provide the the path and extension of a graphics file:
%% \includegraphics{diamondrule} is completely sufficient.
%%
\ifpdf%                                % if we use pdflatex
  \pdfoutput=1\relax                   % create PDFs from pdfLaTeX
  \pdfcompresslevel=9                  % PDF Compression
  \pdfoptionpdfminorversion=7          % create PDF 1.7
  \ExecuteOptions{pdftex}
  \usepackage{graphicx}                % allow us to embed graphics files
  \DeclareGraphicsExtensions{.pdf,.png,.jpg,.jpeg} % for pdflatex we expect .pdf, .png, or .jpg files
\else%                                 % else we use pure latex
  \ExecuteOptions{dvips}
  \usepackage{graphicx}                % allow us to embed graphics files
  \DeclareGraphicsExtensions{.eps}     % for pure latex we expect eps files
\fi%

%% it is recomended to use ``\autoref{sec:bla}'' instead of ``Fig.~\ref{sec:bla}''
\graphicspath{{figures/}{pictures/}{figs/}{./}} % where to search for the images

\usepackage{microtype}                 % use micro-typography (slightly more compact, better to read)
\PassOptionsToPackage{warn}{textcomp}  % to address font issues with \textrightarrow
\usepackage{textcomp}                  % use better special symbols
\usepackage{mathptmx}                  % use matching math font
\usepackage{times}                     % we use Times as the main font
\renewcommand*\ttdefault{txtt}         % a nicer typewriter font
\usepackage{cite}                      % needed to automatically sort the references
\usepackage{tabu}                      % only used for the table example
\usepackage{booktabs}                  % only used for the table example
\usepackage{amsmath,bm}
%% We encourage the use of mathptmx for consistent usage of times font
%% throughout the proceedings. However, if you encounter conflicts
%% with other math-related packages, you may want to disable it.

%% In preprint mode you may define your own headline.
%\preprinttext{To appear in IEEE Transactions on Visualization and Computer Graphics.}

\usepackage[pdftex,dvipsnames]{xcolor}  % Coloured text etc.
\usepackage{listings}

\newcommand{\taoli}[1]{\textcolor{orange}{[#1]}}

\definecolor{codegreen}{rgb}{0,0.6,0}
\definecolor{codegray}{rgb}{0.5,0.5,0.5}
\definecolor{codepurple}{rgb}{0.58,0,0.82}
\definecolor{backcolour}{rgb}{0.95,0.95,0.92}
 
\lstdefinestyle{mystyle}{
    backgroundcolor=\color{backcolour},   
    commentstyle=\color{codegreen},
    keywordstyle=\color{magenta},
    numberstyle=\tiny\color{codegray},
    stringstyle=\color{codepurple},
    basicstyle=\footnotesize,
    breakatwhitespace=false,         
    %breaklines=true,                 
    captionpos=b,                    
    keepspaces=true,                 
    numbers=left,                    
    numbersep=5pt,                  
    showspaces=false,                
    showstringspaces=false,
    showtabs=false,                  
    tabsize=2
}

\lstset{style=mystyle}

\newcommand{\zhimin}[1]{\textcolor{magenta}{[ZHIMIN: #1]}}
\newcommand{\shusen}[1]{\textcolor{red}{[#1]}}

\usepackage{tikz}
\def\checkmark{\tikz\fill[scale=0.4](0,.35) -- (.25,0) -- (1,.7) -- (.25,.15) -- cycle;} 

%% If you are submitting a paper to a conference for review with a double
%% blind reviewing process, please replace the value ``0'' below with your
%% OnlineID. Otherwise, you may safely leave it at ``0''.
\onlineid{1095}

%% declare the category of your paper, only shown in review mode
\vgtccategory{Research}
%% please declare the paper type of your paper to help reviewers, only shown in review mode
%% choices:
%% * algorithm/technique
%% * application/design study
%% * evaluation
%% * system
%% * theory/model
\vgtcpapertype{Application}

%% Paper title.
\title{
Supplementary Materials
%Where and How Do Models Fail? A Visual Exploration Tool for Analyzing and Interpreting Natural Language Inference Models
}
\vgtcinsertpkg

\begin{document}
\firstsection{Appendix A: Detailed description of the Model}
\maketitle

The Decomposable Attention model is an end-to-end model that takes vectorized premise sentence and hypothesis sentence, and output a probability distribution over three labels. The pipeline can be separated into three phases: encoder, attention and classifier. Here we briefly formalize these three layers.

\subsection{Encoder}
Denote the premise sentence as $\overline{\boldsymbol{P}} = (\overline{\boldsymbol{p}}_1, \overline{\boldsymbol{p}}_2, ..., \overline{\boldsymbol{p}}_m)$ and hypothesis sentence as $\overline{\boldsymbol{H}} = (\overline{\boldsymbol{h}}_1, \overline{\boldsymbol{h}}_2, ..., \overline{\boldsymbol{h}}_n)$ where $\overline{\boldsymbol{p}}_i$ and $\overline{\boldsymbol{h}}_j$ are the input embedding vectors respectively. The purpose of encoder is to generate specialized encodings that will participate in the following steps of the model. Define the encoder function as
\begin{align}
  \boldsymbol{P} &= F(\overline{\boldsymbol{P}})\\
  \boldsymbol{H} &= F(\overline{\boldsymbol{H}})
\end{align}

where $F$ function is the encoder layer (e.g. a context-free two-layer ReLU). Let $\boldsymbol{P}$ denote the output premise encoding of size $m \times d$, and $\boldsymbol{H}$ the hypothesis encoding of size $n \times d$ where $d$ is the encoding size.

\subsection{Attention} \label{sec:att}
The job of attention layer is to capture pairwise interaction between encodings of premise and hypothesis. First we compute a similarity score matrix $\boldsymbol{S}$ ($m \times n$) by the following
\begin{align}
  S_{ij} &= \boldsymbol{p}_i ^T \boldsymbol{h}_j
\end{align}

Then attention weights are obtained by normalizing $\boldsymbol{S}$ bidirectionally. We will use $\leftarrow$ and $\rightarrow$ to differentiate the notation. That is, for premise token $\boldsymbol{p}_i$, a probability distribution over all $\boldsymbol{h}_j$`s are done as
\begin{align}
  \overleftarrow{\boldsymbol{a}}_i &= \frac{exp(S_{ij})}{\sum_{j \in h} exp(S_{ij})}
\end{align}

where $\overleftarrow{\boldsymbol{a}}_i$ is of dimension $1 \times n$ naturally. Similarly, for every $\boldsymbol{h}_j$, there is a probability distribution over all premise tokens.
\begin{align}
  \overrightarrow{\boldsymbol{a}}_j &= \frac{exp(S_{ij})}{\sum_{i \in p} exp(S_{ij})} \label{eq:aright}
\end{align}

where $\overrightarrow{\boldsymbol{a}}_j$ is a $1 \times m$ dimension. In case a token is not correlated with any other tokens, the model added an extra artificial token to both premise and hypothesis to allow a word aligning with null.

\subsection{Classifier} \label{sec:cls}
Classifier takes features from previous steps and predict a probability distribution over three labels. First, it aggregates features following the bidirectional attentions as
\begin{align}
  \boldsymbol{u}_i &= G([\overleftarrow{\boldsymbol{a}}_i \cdot \boldsymbol{H} \hspace{1ex} | \hspace{1ex}  \overline{\boldsymbol{p}}_i]) \\
  \boldsymbol{v}_j &= G([\overrightarrow{\boldsymbol{a}}_j \cdot \boldsymbol{P} \hspace{1ex}  | \hspace{1ex}  \overline{\boldsymbol{h}}_j])
\end{align}

where G is the aggregation layer (e.g. two-layer ReLU) and $[\cdot|\cdot]$ denotes concatenation. Notice this aggregation layer jointly considers embeddings, encodings and attentions. With the aggregated premise feature $\boldsymbol{u}_i$ and hypothesis feature $\boldsymbol{v}_j$, the classifier does further feature conjunction to predict weight on each entailment score:
\begin{align}
  \boldsymbol{y} &= Z([\sum_i \boldsymbol{u}_i | \sum_j \boldsymbol{v}_j]) \label{eq:yscore}
\end{align}

where $Z$ is another neural layer (e.g. two-layer ReLU followed by a linear layer). For training, the scores $\boldsymbol{y}$ are connected to multi-class cross entropy loss function. For inference, the predicted label is $\arg\max_i y_i$.

%\taoli{Do we need citation in Appendix?}

\textbf{Comment} We used the same network setting in the Decomposable Attention paper~\cite{parikh2016emnlp}. That are, network hidden size is $200$; input word embeddings are $200$ embeddings (downsampled from $300$ dimensional GloVe embeddings~\cite{PenningtonSocherManning2014} by a learnable linear layer); dropout~\cite{SrivastavaHinton2014} is $0.2$ for learning; and learning rate is $0.05$.

\section{Appendix B: Code Example: Link Visualization with Model}
\begin{lstlisting}[language=Python, caption=Code for generating the visualization.]
from visPackage import *
from modelInterface import *
from sentenceGenerator import *

#initialize NLP model
model = modelInterface(
    wordDict="../data/snli_1.0/snli_1.0.word.dict",
    wordVec="../data/glove.hdf5", 
    model="../data/local_300_parikh")

#sentence perturbation
gen = sentenceGenerator()
#dependency parsing
dep = dependencyTree()

#define visualization components
visLayout = {
    "column":[{"row":["Selector", 
                      "Sentence", 
                      "Pipeline"]},
            {"row":["AttentionGraph", 
                    "AttentionMatrix", 
                    "Prediction"]}]
    }

#create vis interface
vis = textEntailVisModule(visLayout)

#setup model callbacks
vis.setPredictionHook(model.predict)
vis.setAttentionHook(model.attention)
vis.setPredictionUpdateHook(model.updatePrediction)
vis.setAttentionUpdateHook(model.updateAttention)
vis.setReloadModelCallback(model.reloadModel)
vis.setPipelineStatisticHook(model.pipelineStatistic)

#setup perturbation and grammar parsing callbacks
vis.setSentencePerturbationHook(gen.perturbSentence)
vis.setSentenceParserHook(dep.getDependencyTree)

#open browser to show visualization
modelVis.show()

\end{lstlisting}
\bibliographystyle{abbrv-doi}
%\bibliographystyle{abbrv-doi-narrow}
%\bibliographystyle{abbrv-doi-hyperref}
%\bibliographystyle{abbrv-doi-hyperref-narrow}

\bibliography{entailVis.bib}
\end{document}